\documentclass[11pt,a4paper]{article}
\usepackage{amsmath,amssymb,amsthm,bm}
\usepackage{geometry}
\geometry{margin=1in}
\usepackage{hyperref}
\usepackage{graphicx}
\usepackage{setspace}
\usepackage{booktabs}
\onehalfspacing

\title{\textbf{Paper OOL-1: The Thermodynamic Mandate and Dissipative Structuring:\\ A Constraint-Driven Flux Dynamics (CDFD) Perspective}}
\author{Steve Bico Mujjabi, MD\\
\small Independent Researcher\\
\small ORCID: \href{https://orcid.org/0009-0001-0556-5516}{0009-0001-0556-5516}}
\date{May 2026}

\begin{document}
\maketitle

\begin{abstract}
This opening manuscript frames abiogenesis as a nonequilibrium systems problem rather than a single lucky molecular accident. Constraint-Driven Flux Dynamics (CDFD) is used as a coarse-grained language for asking when environmental energy flux, geometric constraint, active surface formation, and memory-like persistence can reinforce one another. The central object is the adaptive ratio
\[
\Psi_s=\left(\frac{\Phi}{C}\right)S M_s,
\]
where $\Phi$ denotes available driving flux, $C$ denotes the effective resistance or constraint field, $S$ denotes responsive organization, and $M_s$ denotes retained structural memory. The paper does not prove that life is inevitable in every driven environment. It establishes the dependency structure that later papers test in geochemical, polymerization, compartment, replication, parasitic, photochemical, and synthesis settings.
\end{abstract}

\section{Introduction}
Living systems maintain local order by coupling themselves to external gradients. The origin-of-life problem is therefore not only the problem of making particular molecules; it is also the problem of sustaining useful work against dilution, hydrolysis, side reactions, and spatial loss. In CDFD language, a prebiotic system becomes interesting when a flux source and a constraint architecture create a persistent opportunity for organization.

This manuscript supplies the series-level vocabulary. This vocabulary is CDFL, the CDFD flow-constraint-memory language used by the public runtime and by the manuscript series. It deliberately keeps the claim at the level of mechanism class: driven systems can favor structures that increase usable dissipation, but only if those structures also retain enough organization to persist.

\section{Scientific Lineage and Complementarity}
Part II stands on a long origin-of-life tradition rather than replacing it. Oparin's coacervate program, Wächtershäuser's surface-metabolism proposal, Russell and Martin's alkaline hydrothermal-vent work, Eigen and Orgel's error-threshold logic, Kauffman's autocatalytic-set theory, Szostak's protocell program, Lane's bioenergetic arguments, Powner's nucleotide chemistry, Patel's cyanosulfidic network chemistry, and Damer's cycling-environment hypothesis all identify real parts of the puzzle~\cite{Wachtershauser1988,Russell1994,Martin2007,Eigen1971,Orgel1963,Kauffman1986,Szostak2010,Lane2010,Powner2009,Patel2015,Damer2020}. CDFD's role is to provide a single dependency grammar across those traditions: what is the input, what constrains it, what routes it, what is retained, and what would falsify the claimed coupling?

This framing is meant to complement those schools, not erase their disagreements. Vent models, hot-spring models, mineral-surface chemistry, RNA-world chemistry, metabolism-first chemistry, and protocell experiments can all be compared by asking whether they solve the same sequence of functional gates.

\textbf{Status: Historical framing}

\section{Relationship to Part I}
Part I of this series introduced CDFL, the CDFD state grammar for flow, constraint, responsiveness, memory, and threshold crossing in fundamental-physics language~\cite{MujjabiPartI2026}. Part II does not require the reader to accept the stronger physical claims of Part I. It uses the shared notation as a disciplined bookkeeping layer: $\Phi$ names drive, $C$ names resistance or constraint, $S$ names responsive organization, and $M_s$ names retained structure. The biological question is whether those variables can be instantiated by geochemical, chemical, and protocellular mechanisms.

\textbf{Status: Framing}


\section{Flux, Constraint, Surface, and Memory}
The shared phenomenological balance used throughout this Part II sequence is
\begin{equation}
    \frac{\partial \Phi}{\partial t}
    =
    \nabla \cdot \left(\frac{1}{C}\nabla\Phi\right)+S-D,
    \label{eq:flux-balance}
\end{equation}
where $D$ collects loss, degradation, washout, and uncoupled dissipation. The terms are interpreted as follows:
\begin{itemize}
    \item \textbf{Flux ($\Phi$)}: chemical, redox, photochemical, thermal, or electrochemical drive available to the system.
    \item \textbf{Constraint ($C$)}: transport resistance, activation barrier, geometric confinement, permeability, or maintenance burden.
    \item \textbf{Responsiveness ($S$)}: the ability of local structure to route flux into productive pathways rather than noise.
    \item \textbf{Systemic memory ($M_s$)}: persistence of structure, sequence, composition, topology, or boundary state across time.
\end{itemize}

All symbols are local phenomenological variables unless explicitly tied to a measurement. They should not be read as universal constants. Later papers redefine specialized symbols only after stating their physical meaning, because the same letter can otherwise mean different things in reaction kinetics, transport theory, and bioenergetics.

The adaptive ratio
\begin{equation}
    \Psi_s=\left(\frac{\Phi}{C}\right)S M_s
    \label{eq:psi}
\end{equation}
is not introduced as a universal constant. It is a bookkeeping device for tracking whether the components needed for persistence are co-present. A high flux without retention burns out. A strong constraint without flux freezes. A responsive surface without memory remains transient.

\section{Dissipative Structuring}
Driven systems can form structures that route energy more effectively than homogeneous diffusion. In prebiotic settings, the relevant question is whether such structures can become chemically consequential: mineral interfaces, pores, films, droplets, polymers, and reaction loops all change the effective relationship between $\Phi$ and $C$.

For a bare chemical soup, one may have $S\approx 0$ and $M_s\to 0$, so transient fluctuations do not accumulate. A successful origin pathway must raise both quantities without disconnecting from the external drive. This is the handoff to Paper 2: the framework needs geological boundary conditions before it can become a concrete origin-of-life argument.

\section{Claim Boundary}
The strongest defensible claim of this paper is conditional:
\begin{quote}
If a prebiotic environment supplies sustained disequilibrium, spatial constraint, and chemistry capable of retaining useful structure, then CDFD predicts selection for configurations that improve coupled throughput and persistence.
\end{quote}

This is not a proof that a particular vent, pool, mineral, or molecule was historically necessary. The later papers narrow those possibilities by adding Fe--S redox chemistry, distributed mineral transport, proton/ion coupling, cycling polymerization, autocatalysis, chemical alphabets, compartments, replication thresholds, parasitic failure modes, and photochemical expansion.

\section{Conclusion}
The thermodynamic opening move is a dependency graph, not a victory lap. Life-like organization requires flux, constraint, responsive routing, and retained memory to reinforce one another. Paper 2 therefore begins the concrete chain with the most immediate physical substrate: iron--sulfur redox and mineral constraint scaffolds.

\section{Reproducibility Note}
This paper is analytic and conceptual. It does not require a notebook or generated output for release. Any toy visualization of Eq.~\eqref{eq:psi} should be treated as an illustration, not as validation of abiogenesis.
\bibliographystyle{plain}
\bibliography{ool_references}
\end{document}
